\section{One genuine unit comparing entire original residue frames}

\begin{definition}\label{def:fc-original-frame-transition-unit}
\lean{EnvelopeProperty.Berkovich.SchemeResidueFrameComparison.transitionUnit}
\leanok
\uses{lem:actual-unit-module-automorphism-scalar}
Let $A$ be any commutative ring and let $P$ be an original module
on $\operatorname{Spec}A$ with two actual entire frames
$e_0,e_1:P\simeq\mathcal O$. The entire original automorphism
$e_1e_0^{-1}$ of the structure module is multiplication by a genuine
global unit. Transport this unit through the actual global-section
isomorphism of $\operatorname{Spec}A$ to obtain a genuine unit
$u(e_0,e_1)\in A^\times$. This choice depends only on the two entire
original frames and precedes all choices of sections and valuations.
\end{definition}
\begin{proof}
The composite of the inverse of $e_0$ with $e_1$ is an actual
structure-module isomorphism. The whole structure-module
classification theorem therefore supplies a global unit whose
actual scalar automorphism equals this entire composite. The
genuine global-section ring isomorphism maps a unit to a unit,
giving the stated original unit in $A$.
\end{proof}

\begin{lemma}\label{lem:fc-original-frame-transition-coefficient}
\lean{EnvelopeProperty.Berkovich.SchemeResidueFrameComparison.transitionUnit_coefficient}
\leanok
\uses{def:fc-original-frame-transition-unit}
For every original whole section $s$ of $P$, its actual coefficients
in the two original frames satisfy $c_1(s)=u(e_0,e_1)c_0(s)$ in $A$.
The same unit works for every original section, including zero.
\end{lemma}
\begin{proof}
Compose the entire scalar-automorphism identity with $e_0$.
The adjacent frame and inverse cancel, leaving $e_1$ equal to
$e_0$ followed by multiplication by the chosen original global
unit. Apply this equality on the original top open to $s$.
The true scalar endomorphism multiplies the original coefficient
by that global unit. Applying the genuine global-section ring
isomorphism preserves multiplication and carries the global unit
to $u(e_0,e_1)$, which proves the exact coefficient identity.
\end{proof}

\begin{lemma}\label{lem:fc-original-frame-transition-positive-value}
\lean{EnvelopeProperty.Berkovich.SchemeResidueFrameComparison.transitionUnit_value_pos,EnvelopeProperty.Berkovich.SchemeResidueFrameComparison.transitionUnit_coefficient_value}
\leanok
\uses{lem:fc-original-frame-transition-coefficient}
For every actual valuation $\nu:A\to\mathbb R_{\geq0}$, the number
$\nu(u(e_0,e_1))$ is strictly positive, and for every original
section one has
$\nu(c_1(s))=\nu(u(e_0,e_1))\nu(c_0(s))$.
\end{lemma}
\begin{proof}
The unit and its inverse multiply to one. Applying the valuation
shows that their values multiply to one. Thus the value of the
unit is nonzero, hence strictly positive in $\mathbb R_{\geq0}$.
Apply the multiplicative valuation to the preceding coefficient
identity to obtain the claimed formula for every section.
\end{proof}

\begin{corollary}\label{cor:fc-native-residue-frame-transition}
\lean{EnvelopeProperty.Berkovich.SchemeResidueFrameComparison.residueFrameCoefficient_change,EnvelopeProperty.Berkovich.SchemeResidueFrameComparison.residueFrameCoefficient_value_change}
\leanok
\uses{lem:fc-original-frame-transition-coefficient,lem:fc-original-frame-transition-positive-value,thm:fc-native-framed-section-residue-transport}
For any actual native valued point $y$ on any original scheme $Y$,
any original module $N$, and any two actual entire frames of its
same original residue fibre, the coefficients of every original
whole section of $N$ differ by this one genuine unit of
$\kappa_Y(y)$. Their actual native valuation values therefore
differ by the one strictly positive value of this unit. No
properness, bounded base, frame-preservation identity, supplied
comparison unit, or norm certificate is required.
\end{corollary}
\begin{proof}
Pull the original whole section through the actual native unit to
the same original residue fibre. The original residue coefficient
is precisely the frame coefficient of this actual pulled section.
Apply the coefficient identity above on the genuine spectrum of
$\kappa_Y(y)$ and then apply the actual native residue valuation.
The positivity statement for the same unit applies to this native
valuation, giving the asserted strictly positive common factor.
\end{proof}
